\section*{Özet}
Prosedürel olarak üretilen çok çözünürlüklü arazilerde, kaba çözünürlüklü
küresel yüzeyin ince detaylı yerel bir yama ile birleştiği sınırda
kullanılan karıştırma (blend) fonksiyonunun şekli, fiziksel olarak
sorgulanabilir temas geometrisinin kalitesini doğrudan etkiler; ancak bu
seçim genellikle salt görsel bir tercih olarak ele alınır ve fiziksel
doğruluk açısından nicel olarak nadiren değerlendirilir. Bu çalışma,
Isaac Lab doğrudan iş akışı ortamında Soft Actor-Critic (SAC) ile
eğitilen, dört bacaklı ve iki eksenli gimbal itkili bir ay iniş aracında
hâlihazırda kullanılan kompakt destekli smoothstep karıştırmasını,
önceden kullanılan sonsuz destekli Gauss karıştırmasıyla aynı geçiş
yarıçapı bütçesi altında karşılaştırmakta ve iniş ayağı temas izinin
yakınında Gauss, ötesinde smoothstep kullanan yeni bir hibrit yöntem
önermektedir.

Otuz tohum üzerindeki analitik ölçümler, üretim öncesi (as-shipped) Gauss
biçiminin smoothstep'e kıyasla yaklaşık 3200 kat daha kötü sınır
yükseklik sürekliliği (C0, RMS) ve 530 kat daha kötü eğim sürekliliği
(C1) ürettiğini (Wilcoxon işaretli sıra testi, $p<\num{1e-8}$) ve bu
kusurun geçiş yarıçapından bağımsız, ölçek değişmez olduğunu
göstermiştir; kalibre edilmiş Gauss ve önerilen hibrit yöntem ise makine
hassasiyetine yakın süreklilik sağlamış, karıştırma şeklinin VRAM
üzerindeki etkisi ise ihmal edilebilir bulunmuştur.

Gerçek Isaac Sim eğitimleri (yöntem başına 5,73 milyon adım) önce tek
tohumla (seed=42) yapılmış ve kalibre edilmiş Gauss ile hibrit'in
smoothstep'e göre belirgin biçimde daha yüksek başarı oranına ulaştığı
görülmüştür. Bu bulgunun güvenilirliğini sınamak için çalışma iki
bağımsız tohumla daha ($n=3$ toplam) genişletilmiş ve yöntemler arasında
hiçbir metrikte (başarı oranı, zaman aşımı oranı, sert iniş oranı)
istatistiksel olarak anlamlı bir fark bulunamamıştır (tüm Mann-Whitney
$p\ge 0{,}20$); ilk ölçümdeki görünür fark büyük ölçüde koşu-koşu (seed)
varyansıyla tutarlıdır. Bu, tek-koşu RL karşılaştırmalarının ne kadar
yanıltıcı olabileceğine dair somut bir örnek sunmakta ve karıştırma
fonksiyonu seçiminin sınır sürekliliği açısından sağlam bir etkiye sahip
olduğunu, ancak bu etkinin nihai RL eğitim performansına -en azından
mevcut ölçek ve $n=3$ ile- ayırt edilebilir biçimde yansımadığını
göstermektedir.

\medskip
\noindent\textbf{Anahtar Kelimeler:} Ay inişi simülasyonu, Takviyeli
öğrenme (Soft Actor-Critic), Isaac Lab, Prosedürel arazi üretimi, Çok
çözünürlüklü karıştırma, Süreklilik analizi
